In a particular ideal refracting telescope of focal ratio $f/5$, the focal
length of the objective lens is 100\,cm and that of the eyepiece is 1\,cm.

\begin{description}
\item[(T7.1)] What is the angular magnification, $m_0$, of the telescope? What
  is the length of the telescope, $L_0$, i.e.\ the distance between its
  objective and eyepiece?
\end{description}

An introduction of a concave lens (Barlow lens) between the objective lens and
the prime focus is a common way to increase the magnification without a large
increase in the length of the telescope. A Barlow lens of focal length 1\,cm is
now introduced between the objective and the eyepiece to double the
magnification.

\begin{description}
\item[(T7.2)] At what distance, $d_{\mathrm{B}}$, from the prime focus must the
  Barlow lens be kept in order to obtain this desired double magnification?
\item[(T7.3)] What is the increase, $\Delta L$, in the length of the telescope?
\end{description}

A telescope is now constructed with the same objective lens and a CCD detector
placed at the prime focus (without any Barlow lens or eyepiece). The size of
each pixel of the CCD detector is 10\,$\mu$m.

\begin{description}
\item[(T7.4)] What will be the distance in pixels between the centroids of the
  images of the two stars, $n_{\mathrm{p}}$, on the CCD, if they are $20''$
  apart on the sky?
\end{description}
